\section{802.11 Device}\label{sec:Device Types / 802.11 Device}

% TODO: normative statements
% TODO: add non-normative paragraph that explains how both FullMAC and SoftMAC can work with this device spec.

The 802.11 device supports the transmission and reception of \hyperref[intro:IEEE 802.11]{802.11} frames.
It includes additional functionality to support Station Mode workflows such as authentication.

\subsection{Device ID}\label{sec:Device Types / 802.11 Device / Device ID}

10

\subsection{Virtqueues}\label{sec:Device Types / 802.11 Device / Virtqueues}

\begin{description}
\item[0] controlq
\item[1] eventq
\item[2] txq
\item[3] rxq
\end{description}

The driver enqueues control requests to the controlq to change or query configuration.

The device enqueues asynchronous events (such as scan results) to the eventq.

TODO: normative statement, driver MUST populate rxq and eventq with appropriately large buffers on initialization.

\subsection{Feature bits}\label{sec:Device Types / 802.11 Device / Feature bits}

None.

TODO or...:

\begin{description}
\item[VIRTIO_WIFI_F_AUTHENTICATE (0)] Device supports authentication commands. % TODO which?
\item[VIRTIO_WIFI_F_CONNECT (1)] Device supports connection commands. % TODO which?
\item[VIRTIO_WIFI_F_CHANGE_INTERFACE (2)] Device supports VIRTIO_WIFI_CMD_CHANGE_INTERFACE.
\item[VIRTIO_WIFI_F_START_AP (3)] Device supports VIRTIO_WIFI_CMD_START_AP.
\item[VIRTIO_WIFI_F_REGISTER_FRAME (4)] Device supports VIRTIO_WIFI_CMD_REGISTER_FRAME.
\item[VIRTIO_WIFI_F_AP_SME (5)] % TODO
\end{description}

\subsection{Device configuration layout}\label{sec:Device Types / 802.11 Device / Device configuration layout}

\begin{lstlisting}
struct virtio_wifi_config {
	__le32 num_of_ifaces;
	__le32 flags;
	__le32 iftypes;
	__le32 available_antennas_tx;
	__le32 available_antennas_rx;
	__le32 n_cipher_suites;
} __attribute__((packed));
\end{lstlisting}

\subsection{Device Initialization}\label{sec:Device Types / 802.11 Device / Device Initialization}

The driver populates rxq and eventq with appropriately large buffers.

The driver determines the number of existing interfaces, and queries them with VIRTIO_WIFI_CMD_GET_INTERFACES.

\subsection{Device Operation}\label{sec:Device Types / 802.11 Device / Device Operation}

The driver makes a request available in the controlq. The device fills
in the response and uses the buffer. The controlq uses common request
and response headers.

The device can send asynchronous events to the driver via the eventq.

The driver can transmit packets by placing them to the txq.

The device places received packets in the rxq.

For the control, rx and tx queues a common header is defined for requests and responses:

\begin{lstlisting}
/* common request and response header */
struct virtio_wifi_hdr {
	__le32 code;
	__le32 index;
};
\end{lstlisting}

\begin{description}
\item[\field{code}] identifies the command type.
\item[\field{index}] indicates what phy or interface index a
command refers to.
\end{description}

In a response, if \field{code} is not VIRTIO_WIFI_S_OK, the value of other response
fields is undefined.

The following commands are defined for the controlq:

\begin{lstlisting}
enum {
	VIRTIO_WIFI_CMD_INVALID = 0,
	VIRTIO_WIFI_CMD_GET_INTERFACES,
	VIRTIO_WIFI_CMD_NEW_INTERFACE,
	VIRTIO_WIFI_CMD_CHANGE_INTERFACE,
	VIRTIO_WIFI_CMD_DELETE_INTERFACE,
	VIRTIO_WIFI_CMD_GET_BAND,
	VIRTIO_WIFI_CMD_GET_BAND_CHANNEL,
	VIRTIO_WIFI_CMD_GET_BAND_BITRATE,
	VIRTIO_WIFI_CMD_GET_CIPHERSUITES,
	VIRTIO_WIFI_CMD_TRIGGER_SCAN,
	VIRTIO_WIFI_CMD_CONNECT,
	VIRTIO_WIFI_CMD_AUTHENTICATE,
	VIRTIO_WIFI_CMD_ASSOCIATE,
	VIRTIO_WIFI_CMD_DEAUTHENTICATE,
	VIRTIO_WIFI_CMD_DISASSOCIATE,
	VIRTIO_WIFI_CMD_ADD_KEY,
	VIRTIO_WIFI_CMD_SET_KEY,
	VIRTIO_WIFI_CMD_GET_KEY,
	VIRTIO_WIFI_CMD_DEL_KEY,
	VIRTIO_WIFI_CMD_SET_STATION,
	VIRTIO_WIFI_CMD_IF_OPERSTATE,
	VIRTIO_WIFI_CMD_GET_SURVEY,
	VIRTIO_WIFI_CMD_START_AP,
	VIRTIO_WIFI_CMD_REGISTER_FRAME,

	/* common status codes */
	VIRTIO_WIFI_S_OK = 0x8000,
	VIRTIO_WIFI_S_BAD_MSG,
	VIRTIO_WIFI_S_NOT_SUPP,
	VIRTIO_WIFI_S_IO_ERR,
	VIRTIO_WIFI_S_NET_DOWN,
	VIRTIO_WIFI_S_NOT_CONN,
	VIRTIO_WIFI_S_NO_LINK
};
\end{lstlisting}

The following commands are defined for the txq and rxq:

\begin{lstlisting}
/* rxq/txq commands */
enum {
	VIRTIO_WIFI_802_3,
	VIRTIO_WIFI_802_11,
};
\end{lstlisting}

\drivernormative{\subsubsection}{Device Operation}{Device Types / 802.11 Device / Device Operation}

The driver MUST interpret response fields other than the \field{struct
virtio_wifi_hdr} field \field{code} only when \field{code} is
VIRTIO_WIFI_S_OK.

The driver MUST allocate enough space for the response in a
device-writable controlq buffer.

\devicenormative{\subsubsection}{Device Operation}{Device Types / 802.11 Device / Device Operation}

For \field{struct virtio_wifi_hdr}, the device MUST set the
\field{code} field to VIRTIO_WIFI_S_OK if the device successfully
executed the request.

For \field{struct virtio_wifi_hdr}, the device MUST set the
\field{code} field to a status other than VIRTIO_WIFI_S_OK if the
device did not successfully execute the request.

For \field{struct virtio_wifi_hdr}, the device MUST set the
\field{code} field to VIRTIO_WIFI_S_EOPNOTSUPP if the device could not
execute the specific request due to an implementation limitation.

For \field{struct virtio_wifi_hdr}, the device MUST set the
\field{code} field to VIRTIO_WIFI_S_EOPNOTSUPP for a request with a
value of the request \field{code} field which is not described in this
specification.

For \field{struct virtio_wifi_hdr}, the device MUST set the
\field{code} field to VIRTIO_WIFI_S_ENODEV if the \field{index}
field value supplied with the request does not identify a phy or an interface.

For \field{struct virtio_wifi_hdr}, the device MUST set the
\field{code} field to VIRTIO_WIFI_S_EINVAL if the request values are
inconsistent with the specification and if the inconsistence is not
described by the requirements which stipulate status
VIRTIO_WIFI_S_EOPNOTSUPP or VIRTIO_WIFI_S_ENODEV.

For \field{struct virtio_wifi_hdr}, the device MUST set the
\field{code} field to VIRTIO_WIFI_S_EINVAL if the request specified in
the request header through the \field{code} field does not fit into
the device read-only buffer.

For \field{struct virtio_wifi_hdr}, the device MUST set the
\field{code} field to VIRTIO_WIFI_S_EINVAL if the response specified in
the request header through the \field{code} field does not fit into
the device write-only buffer, unless the \field{code} field does not
fit into the device write-only buffer.

For \field{struct virtio_wifi_hdr}, the device MUST NOT set the
\field{code} field if the \field{code} field does not fit into the
device write-only buffer.

If the device read-only buffer is bigger than the size of the request
specified in the request header, the device MUST ignore the additional
space.

If the device write-only buffer is bigger than the size of the response
corresponding to the request header, the device MUST ignore the
additional space.

The device MUST set \emph{reserved} fields in a device-writable buffer
to zero.

The device MUST NOT set bits in \emph{flags} fields in a device-writable
buffer which are not allowed according to the negotiated features.

\subsubsection{Common Definitions}\label{sec:Device Types / 802.11 Device / Device Operation / Common Definitions}

This section makes common definitions.

\paragraph{Type-Length-Value Header}\label{sec:Device Types / 802.11 Device / Device Operation / Common Definitions / Type-Length-Value Header}

Many requests and responses carry data fields of variable size.

These data fields are placed in a queue after the request or response structure.

Each data field consists of a Type-Length header followed by the content of the field.

\begin{lstlisting}
/* common type-length-value header */
struct virtio_wifi_tlv {
	__le32 type;
	__le32 length;
	__le32 value[];
};
\end{lstlisting}

\paragraph{Authentication Types}\label{sec:Device Types / 802.11 Device / Device Operation / Common Definitions / Authentication Types}

\begin{lstlisting}
enum virtio_wifi_auth_type {
	VIRTIO_WIFI_AUTHTYPE_OPEN_SYSTEM = 0,
	VIRTIO_WIFI_AUTHTYPE_SHARED_KEY,
	VIRTIO_WIFI_AUTHTYPE_FT,
	VIRTIO_WIFI_AUTHTYPE_NETWORK_EAP,
	VIRTIO_WIFI_AUTHTYPE_SAE,
	VIRTIO_WIFI_AUTHTYPE_FILS_SK,
	VIRTIO_WIFI_AUTHTYPE_FILS_SK_PFS,
	VIRTIO_WIFI_AUTHTYPE_FILS_PK,
	VIRTIO_WIFI_AUTHTYPE_EPPKE,
};
\end{lstlisting}

\paragraph{Interface Types}\label{sec:Device Types / 802.11 Device / Device Operation / Common Definitions / Interface Types}

The following interface types are defined:

\begin{lstlisting}
enum virtio_wifi_iftype {
	VIRTIO_WIFI_IFTYPE_UNSPECIFIED = 0,
	VIRTIO_WIFI_IFTYPE_ADHOC,
	VIRTIO_WIFI_IFTYPE_STATION,
	VIRTIO_WIFI_IFTYPE_AP,
	VIRTIO_WIFI_IFTYPE_AP_VLAN,
	VIRTIO_WIFI_IFTYPE_WDS,
	VIRTIO_WIFI_IFTYPE_MONITOR,
	VIRTIO_WIFI_IFTYPE_MESH_POINT,
	VIRTIO_WIFI_IFTYPE_P2P_CLIENT,
	VIRTIO_WIFI_IFTYPE_P2P_GO,
	VIRTIO_WIFI_IFTYPE_P2P_DEVICE,
	VIRTIO_WIFI_IFTYPE_OCB,
	VIRTIO_WIFI_IFTYPE_NAN,
	VIRTIO_WIFI_IFTYPE_MAX,
};
\end{lstlisting}

The following definition dictates the address byte length:

\begin{lstlisting}
#define ETH_ALEN 6 /* Octets in one ethernet addr	 */
\end{lstlisting}

The following definition dictates the maximum allowed SSID length:

\begin{lstlisting}
#define IEEE80211_MAX_SSID_LEN 32
\end{lstlisting}

The following definitions dictate the maximum allowed cipher suites and authentication and key management (AKM) suites:

\begin{lstlisting}
#define MAX_NR_CIPHER_SUITES 5
#define MAX_NUM_AKM_SUITES 10
\end{lstlisting}

\paragraph{Deauthentication And Disassociation types}\label{sec:Device Types / 802.11 Device / Device Operation / Common Definitions / Deauthentication And Disassociation types}

\begin{lstlisting}
/* Used both for deauth and disassoc */
struct virtio_wifi_disconnect_request {
	struct virtio_wifi_hdr hdr;
	__u8 mac[ETH_ALEN];
	__le16 reason_code;
	__u8 local_state_change;
	__u8 padding[7];
};

enum virtio_wifi_disconnect_request_field {
	VIRTIO_WIFI_DISCONN_REQ_IE = 0,
};
\end{lstlisting}

\begin{description}
\item[\field{mac}] TODO
\item[\field{reason_code}] TODO
\item[\field{local_state_change}] TODO
\item[\field{VIRTIO_WIFI_DISCONN_REQ_IE}] contains a request Information Element (IE) and is optional.
\end{description}

\subsubsection{Events}\label{sec:sec:Device Types / 802.11 Device / Device Operation / Events}

\begin{lstlisting}
enum virtio_wifi_event_code {
	VIRTIO_WIFI_EVENT_CODE_INVALID = 0,
	VIRTIO_WIFI_EVENT_CODE_SCAN_RESULT,
	VIRTIO_WIFI_EVENT_CODE_END_SCAN,
	VIRTIO_WIFI_EVENT_CODE_CANCEL_SCAN,
	VIRTIO_WIFI_EVENT_CODE_AUTHENTICATE,
	VIRTIO_WIFI_EVENT_CODE_AUTHENTICATE_TIMED_OUT,
	VIRTIO_WIFI_EVENT_CODE_ASSOCIATE,
	VIRTIO_WIFI_EVENT_CODE_ASSOCIATE_TIMED_OUT,
	VIRTIO_WIFI_EVENT_CODE_CONNECT,
	VIRTIO_WIFI_EVENT_CODE_CONNECT_TIMED_OUT,
	VIRTIO_WIFI_EVENT_CODE_DISCONNECT,
	VIRTIO_WIFI_EVENT_CODE_DISASSOCIATE,
	VIRTIO_WIFI_EVENT_CODE_DEAUTHENTICATE,
	VIRTIO_WIFI_EVENT_CODE_NEW_STATION,
	VIRTIO_WIFI_EVENT_CODE_DELETE_STATION,
};

struct virtio_wifi_event {
	__le32 code;
	__le32 ifindex;
};
\end{lstlisting}

Each \field{struct virtio_wifi_event} might be embedded in another structure, depending on the event code type.

\subsubsection{Querying Physical Properties}\label{sec:Device Types / 802.11 Device / Device Operation / Querying Physical Properties}

The driver can query physical properties of the device. These do not change during the lifetime of a device.

% TODO normative statement about physical properties not changing.

\paragraph{Frequency Bands}\label{sec:Device Types / 802.11 Device / Device Operation / Querying Physical Properties / Frequency Bands}

VIRTIO_WIFI_CMD_GET_BAND

TODO: explain

\begin{lstlisting}
/**
 * enum virtio_wifi_freq_band - Frequency band
 * @VIRTIO_WIFI_BAND_2GHZ: 2.4 GHz ISM band
 * @VIRTIO_WIFI_BAND_5GHZ: around 5 GHz band (4.9 - 5.7 GHz)
 * @VIRTIO_WIFI_BAND_60GHZ: around 60 GHz band (58.32 - 69.12 GHz)
 * @VIRTIO_WIFI_BAND_6GHZ: around 6 GHz band (5.9 - 7.2 GHz)
 * @VIRTIO_WIFI_BAND_S1GHZ: around 900MHz, supported by S1G PHYs
 * @VIRTIO_WIFI_BAND_LC: light communication band (placeholder)
 */
enum virtio_wifi_freq_band {
	VIRTIO_WIFI_FREQ_BAND_2GHZ = 0,
	VIRTIO_WIFI_FREQ_BAND_5GHZ,
	VIRTIO_WIFI_FREQ_BAND_60GHZ,
	VIRTIO_WIFI_FREQ_BAND_6GHZ,
	VIRTIO_WIFI_FREQ_BAND_S1GHZ,
	VIRTIO_WIFI_FREQ_BAND_LC,
};

/* enum ieee80211_edmg_bw_config */
enum virtio_wifi_edmg_bw_config {
	VIRTIO_WIFI_EDMG_BW_CONFIG_INVALID = 0,
	VIRTIO_WIFI_EDMG_BW_CONFIG_4 = 4,
	VIRTIO_WIFI_EDMG_BW_CONFIG_5 = 5,
	VIRTIO_WIFI_EDMG_BW_CONFIG_6 = 6,
	VIRTIO_WIFI_EDMG_BW_CONFIG_7 = 7,
	VIRTIO_WIFI_EDMG_BW_CONFIG_8 = 8,
	VIRTIO_WIFI_EDMG_BW_CONFIG_9 = 9,
	VIRTIO_WIFI_EDMG_BW_CONFIG_10 = 10,
	VIRTIO_WIFI_EDMG_BW_CONFIG_11 = 11,
	VIRTIO_WIFI_EDMG_BW_CONFIG_12 = 12,
	VIRTIO_WIFI_EDMG_BW_CONFIG_13 = 13,
	VIRTIO_WIFI_EDMG_BW_CONFIG_14 = 14,
	VIRTIO_WIFI_EDMG_BW_CONFIG_15 = 15,
};

struct virtio_wifi_band {
  struct virtio_wifi_hdr hdr;
  /* enum virtio_wifi_freq_band */
  __le32 band;
  __le32 n_channels;
  __le32 n_bitrates;
  __le16 n_iftype_data;
  struct virtio_wifi_sta_ht_cap {
    __le16 cap; /* use IEEE80211_HT_CAP_ */
    __u8 ht_supported;
    __u8 ampdu_factor;
    __u8 ampdu_density;
    __u8 padding;
    struct mcs_info {
      __le16 rx_highest;
      #define VIRTIO_WIFI_HT_MCS_MASK_LEN 10
      __u8 rx_mask[VIRTIO_WIFI_HT_MCS_MASK_LEN];
      __u8 tx_params;
      __u8 reserved;
    } mcs;
  } ht_cap;
  __le16 padding;
  struct sta_vht_cap {
    __le32 vht_supported;
    __le32 cap; /* use IEEE80211_VHT_CAP_ */
    struct virtio_wifi_vht_mcs_info {
      __le16 rx_mcs_map;
      __le16 rx_highest;
      __le16 tx_mcs_map;
      __le16 tx_highest;
    } vht_mcs;
  } vht_cap;
  /* struct ieee80211_sta_s1g_cap */
  struct sta_s1g_cap {
    __u8 s1g;
    __u8 cap[10]; /* use S1G_CAPAB_ */
    __u8 nss_mcs[5];
  } s1g_cap;
  /* struct ieee80211_edmg */
  struct virtio_wifi_edmg {
    __u8 channels;
    __u8 padding[3];
    /* enum virtio_wifi_edmg_bw_config */
    __le32 bw_config;
  } edmg_cap;
  };
\end{lstlisting}

\paragraph{Channels}\label{sec:Device Types / 802.11 Device / Device Operation / Querying Physical Properties / Channels}

VIRTIO_WIFI_CMD_GET_BAND_CHANNEL
TODO: explain

\begin{lstlisting}
struct virtio_wifi_cmd_get_band_channel_request {
	/* VIRTIO_WIFI_CMD_GET_BAND_CHANNEL */
	struct virtio_wifi_hdr hdr;
	__le32 channel_index;
};
\end{lstlisting}

\begin{lstlisting}
enum virtio_wifi_chan_width {
	VIRTIO_WIFI_CHAN_WIDTH_20_NOHT = 0,
	VIRTIO_WIFI_CHAN_WIDTH_20,
	VIRTIO_WIFI_CHAN_WIDTH_40,
	VIRTIO_WIFI_CHAN_WIDTH_80,
	VIRTIO_WIFI_CHAN_WIDTH_80P80,
	VIRTIO_WIFI_CHAN_WIDTH_160,
	VIRTIO_WIFI_CHAN_WIDTH_5,
	VIRTIO_WIFI_CHAN_WIDTH_10,
	VIRTIO_WIFI_CHAN_WIDTH_1,
	VIRTIO_WIFI_CHAN_WIDTH_2,
	VIRTIO_WIFI_CHAN_WIDTH_4,
	VIRTIO_WIFI_CHAN_WIDTH_8,
	VIRTIO_WIFI_CHAN_WIDTH_16,
	VIRTIO_WIFI_CHAN_WIDTH_320,
};

struct virtio_wifi_channel {
	struct virtio_wifi_hdr hdr;
	/* enum virtio_wifi_freq_band */
	__le32 band;
	__le32 center_freq;
	__le32 center_freq1;
	__le32 center_freq2;
	__le16 freq_offset;
	__le16 hw_value;
	/* enum virtio_wifi_chan_width */
	__le32 width;
	__le32 flags;
	__le32 max_antenna_gain;
	__le32 max_power;
	__le32 max_reg_power;
	__le32 dfs_state_entered;
	__le32 dfs_cac_ms;
	__u8 beacon_found;
	__u8 psd;
	__le16 reserved;
};
\end{lstlisting}

\paragraph{Bitrates}\label{sec:Device Types / 802.11 Device / Device Operation / Querying Physical Properties / Bitrates}

VIRTIO_WIFI_CMD_GET_BAND_BITRATE

\begin{lstlisting}
struct virtio_wifi_cmd_get_band_bitrate_request {
	/* VIRTIO_WIFI_CMD_GET_BAND_BITRATE */
	struct virtio_wifi_hdr hdr;
	__le32 band_index;
};
\end{lstlisting}

\begin{lstlisting}
struct virtio_wifi_band_bitrate {
	struct virtio_wifi_hdr hdr;
	__le32 flags;
	__le16 bitrate;
	__le16 hw_value;
	__le16 hw_value_short;
	__le16 reserved;
};
\end{lstlisting}

TODO explain get band bitrate.

\subsubsection{Interface Management}\label{sec:Device Types / 802.11 Device / Device Operation / Interface Management}

\paragraph{Querying Interfaces}\label{sec:Device Types / 802.11 Device / Device Operation / Interface Management / Querying Interfaces}

The driver can retrieve information about existing interfaces with the VIRTIO_WIFI_CMD_GET_INTERFACES command.

The device responds with an array of \field{struct virtio_wifi_interface} values of length \field{num_of_ifaces}.

\paragraph{Creating Interfaces}\label{sec:Device Types / 802.11 Device / Device Operation / Interface Management / Creating Interfaces}

To create an interface, the driver submits a \field{struct virtio_wifi_interface} to the controlq with \field{code} set to VIRTIO_WIFI_CMD_NEW_INTERFACE and \field{index} set to the appropriate wiphy index.

The device responds with a \field{struct virtio_wifi_interface} with updated fields indicating success or a \field{struct virtio_wifi_hdr} on failure.

\begin{lstlisting}
enum virtio_wifi_iftype {
	VIRTIO_WIFI_IFTYPE_UNSPECIFIED = 0,
	VIRTIO_WIFI_IFTYPE_ADHOC,
	VIRTIO_WIFI_IFTYPE_STATION,
	VIRTIO_WIFI_IFTYPE_AP,
	VIRTIO_WIFI_IFTYPE_AP_VLAN,
	VIRTIO_WIFI_IFTYPE_WDS,
	VIRTIO_WIFI_IFTYPE_MONITOR,
	VIRTIO_WIFI_IFTYPE_MESH_POINT,
	VIRTIO_WIFI_IFTYPE_P2P_CLIENT,
	VIRTIO_WIFI_IFTYPE_P2P_GO,
	VIRTIO_WIFI_IFTYPE_P2P_DEVICE,
	VIRTIO_WIFI_IFTYPE_OCB,
	VIRTIO_WIFI_IFTYPE_NAN,
	VIRTIO_WIFI_IFTYPE_MAX,
};

enum virtio_wifi_interface_flags {
	VIRTIO_WIFI_INTERFACE_F_ASSOC = 0,
	VIRTIO_WIFI_INTERFACE_F_IBSS_JOINED,
	VIRTIO_WIFI_INTERFACE_F_IBSS_CREATOR,
	VIRTIO_WIFI_INTERFACE_F_PS,
	VIRTIO_WIFI_INTERFACE_F_S1G,
	VIRTIO_WIFI_INTERFACE_F_IDLE,
};

struct virtio_wifi_interface {
	struct virtio_wifi_hdr hdr;
	/* enum virtio_wifi_iftype */
	__le32 iftype;
	__le32 index;
  __u8 ifname[16];
	/* struct ieee80211_vif_cfg */
	struct virtio_wifi_interface_cfg {
		/* VIRTIO_WIFI_INTERFACE_F_X */
		__le32 flags;
		__le16 aid;
		__le16 eml_cap;
		__le16 eml_med_sync_delay;
		__le16 mld_capa_op;
#define IEEE80211_BSS_ARP_ADDR_LIST_LEN 4
		__be32 arp_addr_list[IEEE80211_BSS_ARP_ADDR_LIST_LEN];
		__le32 arp_addr_cnt;
		__u8 ssid[IEEE80211_MAX_SSID_LEN];
		__le16 ssid_len;
		__u8 ap_addr[ETH_ALEN];
	} cfg;
	__u8 addr[ETH_ALEN];
	__le16 reserved;
	__le32 driver_flags;
};
\end{lstlisting}

\paragraph{Deleting Interfaces}\label{sec:Device Types / 802.11 Device / Device Operation / Interface Management / Deleting Interfaces}

To delete an interface, the driver submits a \field{struct virtio_wifi_hdr} to the controlq with \field{code} set to VIRTIO_WIFI_CMD_DELETE_INTERFACE and \field{index} set to the appropriate interface index.

The device responds with a \field{struct virtio_wifi_hdr} indicating success or failure.

\paragraph{Setting Operational State}\label{sec:Device Types / 802.11 Device / Device Operation / Interface Management / Setting Operational State}

% TODO

\begin{lstlisting}
enum virtio_wifi_if_operstate {
	VIRTIO_WIFI_IF_OPER_UNKNOWN = 0,
	VIRTIO_WIFI_IF_OPER_NOTPRESENT = 1,
	VIRTIO_WIFI_IF_OPER_DOWN = 2,
	VIRTIO_WIFI_IF_OPER_LOWERLAYERDOWN = 3,
	VIRTIO_WIFI_IF_OPER_TESTING = 4,
	VIRTIO_WIFI_IF_OPER_DORMANT = 5,
	VIRTIO_WIFI_IF_OPER_UP = 6,
};

struct virtio_wifi_set_if_operstate {
	struct virtio_wifi_hdr hdr;
	__le32 flags;
	__le32 change_mask;
};
\end{lstlisting}

\subsubsection{Triggering A Scan}\label{sec:Device Types / 802.11 Device / Device Operation / Scan}

Scanning is an asynchronous operation: the driver requests a scan, and the device replies with VIRTIO_WIFI_EVENT_CODE_SCAN_RESULT events each with a scan result, and concludes the scan with VIRTIO_WIFI_EVENT_CODE_END_SCAN event for a successful scan or VIRTIO_WIFI_EVENT_CODE_CANCEL_SCAN to signify failure.

The driver begins a scan by sending a \field{struct virtio_wifi_hdr} with \field{code} set to  VIRTIO_WIFI_CMD_TRIGGER_SCAN and \field{index} to the interface that wants to perform the scan. The device responds with VIRTIO_WIFI_S_OK if it has accepted the operation.

For each scan result, the device notifies the driver with an asynchronous event with \field{code} as VIRTIO_WIFI_EVENT_CODE_SCAN_RESULT and \field{index} set to the interface index.

The full event type is:

%TODO convert to TLV

\begin{lstlisting}
struct virtio_wifi_scan_result {
	struct virtio_wifi_event hdr;
	/* the TSF sent by the peer in the beacon/probe response (or 0) */
	__le64 tsf;
	/* frame type (if known) enum cfg80211_bss_frame_type */
	__le32 ftype;
	/* the capability field sent by the peer */
	__le16 capability;
	/* the beacon interval announced by the peer */
	__le16 beacon_interval;
	__le32 band;
	__le32 center_freq;
	__le32 freq_offset;
	__le16 hw_value;
	__u8 bssid[ETH_ALEN];
	__le32 signal;
};

enum virtio_wifi_scan_result_field {
       VIRTIO_WIFI_SCAN_RESULT_IE = 0,
};
\end{lstlisting}

\begin{description}
\item[\field{VIRTIO_WIFI_SCAN_RESULT_IE}] contains the scan result's Information Element (IE).
\end{description}

\subsubsection{Association}\label{sec:Device Types / 802.11 Device / Device Operation / Association}

The driver sends a \field{struct virtio_wifi_association_request} with the header \field{code} set to VIRTIO_WIFI_CMD_ASSOCIATION and the header \field{index} set to the interface index.

The device responds with VIRTIO_WIFI_S_OK if it has accepted the operation.

\begin{lstlisting}
struct virtio_wifi_association_request {
	struct virtio_wifi_event hdr;
	__u8 mac[ETH_ALEN];
	__le16 ssid_len;
	__u8 ssid[IEEE80211_MAX_SSID_LEN];
	__le32 freq;
	__le32 freq_offset;
	__le32 wpa_versions;
	__le32 cipher_group;
	__le32 n_ciphers_pairwise;
	__le32 ciphers_pairwise[MAX_NR_CIPHER_SUITES];
	__le32 n_akm_suites;
	__le32 akm_suites[MAX_NUM_AKM_SUITES];
	__le32 padding0;
};

enum virtio_wifi_association_request_field {
       VIRTIO_WIFI_ASSOC_REQ_IE = 0,
};
\end{lstlisting}

\begin{description}
\item[\field{VIRTIO_WIFI_ASSOC_REQ_IE}] contains the association requests's Information Element (IE) and is optional.
\end{description}

The device responds with asynchronous VIRTIO_WIFI_EVENT_CODE_ASSOCIATE or
	VIRTIO_WIFI_EVENT_CODE_ASSOCIATE_TIMED_OUT events.

The VIRTIO_WIFI_EVENT_CODE_ASSOCIATE event has the type:

%TODO convert to TLVs

\begin{lstlisting}
struct virtio_wifi_association_response {
	struct virtio_wifi_event hdr;
	__le64 frame_len;
	__le64 req_ies_len;
	__le32 uapsd_queues;
	__le32 padding0;
};

enum virtio_wifi_association_response_field {
       VIRTIO_WIFI_ASSOC_RESP_FRAME = 0,
       VIRTIO_WIFI_ASSOC_RESP_IES,
};
\end{lstlisting}

\begin{description}
  \item[\field{VIRTIO_WIFI_ASSOC_RESP_FRAME}] contains the association response's 802.11 frame and is optional.

\item[\field{VIRTIO_WIFI_ASSOC_RESP_IES}] contains the association response's Information Element (IE) and is optional.
\end{description}

To disassociate, the driver sends a \field{struct virtio_wifi_disconnect_request} with \field{code} set to VIRTIO_WIFI_CMD_DISASSOCIATE.

The device sends an asynchronous VIRTIO_WIFI_EVENT_CODE_DISASSOCIATE event whenever a disassociation happens.

\subsubsection{Authentication}\label{sec:Device Types / 802.11 Device / Device Operation / Authentication}

The driver sends a \field{struct virtio_wifi_authentication_request} with the header \field{code} set to VIRTIO_WIFI_CMD_AUTHENTICATE and the header \field{index} set to the interface index.

The device responds with VIRTIO_WIFI_S_OK if it has accepted the operation.

\begin{lstlisting}
struct virtio_wifi_authentication_request {
	struct virtio_wifi_event hdr;
	__u8 mac[ETH_ALEN];
	__le16 ssid_len;
	__u8 ssid[IEEE80211_MAX_SSID_LEN];
	__le32 freq;
	__le32 freq_offset;
	__le32 wpa_versions;
	__le32 cipher_group;
	__le32 n_ciphers_pairwise;
	__le32 ciphers_pairwise[MAX_NR_CIPHER_SUITES];
	__le32 n_akm_suites;
	__le32 akm_suites[MAX_NUM_AKM_SUITES];
	__u8 privacy;
	__u8 padding[3];
};

enum virtio_wifi_authentication_request_field {
       VIRTIO_WIFI_AUTH_REQ_AUTH_TYPE = 0,
       VIRTIO_WIFI_AUTH_REQ_IE,
};
\end{lstlisting}

\begin{description}
\item[\field{VIRTIO_WIFI_AUTH_AUTH_TYPE}] contains the authentication type of type \field{enum virtio_wifi_auth_type} and is optional.

If omitted, it is assummed it is assigned automatically.

\item[\field{VIRTIO_WIFI_AUTH_REQ_IE}] contains the authentication requests's Information Element (IE) and is optional.
\end{description}

The device responds with asynchronous VIRTIO_WIFI_EVENT_CODE_AUTHENTICATE or
	VIRTIO_WIFI_EVENT_CODE_AUTHENTICATE_TIMED_OUT events.

To deauthenticate, the driver sends a \field{struct virtio_wifi_disconnect_request} with \field{code} set to VIRTIO_WIFI_CMD_DEAUTHENTICATE.

The device sends an asynchronous VIRTIO_WIFI_EVENT_CODE_DEAUTHENTICATE event whenever a deauthentication happens.

\subsubsection{Connection}\label{sec:Device Types / 802.11 Device / Device Operation / Connection}

The device notifies the driver with asynchronous VIRTIO_WIFI_EVENT_CODE_CONNECT or
	VIRTIO_WIFI_EVENT_CODE_DISCONNECT events when a connection change happens.

\begin{lstlisting}
struct virtio_wifi_connection_response {
	struct virtio_wifi_event hdr;
	__u8 mac[ETH_ALEN];
	__le16 status_code;
};

enum virtio_wifi_connection_response_field {
       VIRTIO_WIFI_CONN_REQ_IE = 0,
       VIRTIO_WIFI_CONN_RESP_IE,
};
\end{lstlisting}

\begin{description}
\item[\field{VIRTIO_WIFI_CONN_REQ_IE}] contains the connection requests's Information Element (IE) and is optional.
\item[\field{VIRTIO_WIFI_CONN_RESP_IE}] contains the connection response's Information Element (IE) and is optional.
\end{description}

\begin{lstlisting}
struct virtio_wifi_disconnection_response {
	struct virtio_wifi_event hdr;
	__le16 reason_code;
	__u8 disconnected_by_ap;
  __u8 padding[5];
  __le64 ie_len
};

enum virtio_wifi_disconnection_response_field {
       VIRTIO_WIFI_DISCONN_RESP_IE = 0,
};
\end{lstlisting}

\begin{description}
\item[\field{VIRTIO_WIFI_DISCONN_RESP_IE}] contains the disconnection response's Information Element (IE) and is optional.
\end{description}

\subsubsection{Starting an Access Point}\label{sec:Device Types / 802.11 Device / Device Operation / Starting an Access Point}

% TODO VIRTIO_WIFI_CMD_START_AP

\begin{lstlisting}
struct virtio_wifi_beacon {
	struct virtio_wifi_hdr hdr;
	struct virtio_wifi_channel channel;
	__le32 beacon_interval;
	__le32 beacon_rate;
	__le32 wpa_versions;
	/* enum virtio_wifi_freq_band */
	__le32 band;
	__le32 dtim_period;
	__le32 cipher_group;
	__le32 n_ciphers_pairwise;
	__le32 ciphers_pairwise[MAX_NR_CIPHER_SUITES];
	__le32 n_akm_suites;
	__le32 akm_suites[MAX_NUM_AKM_SUITES];
	__u8 privacy;
	__u8 padding[3];
};

enum virtio_wifi_beacon_field {
	VIRTIO_WIFI_BEACON_HEAD = 0,
	VIRTIO_WIFI_BEACON_TAIL,
	VIRTIO_WIFI_BEACON_IES,
	VIRTIO_WIFI_BEACON_IES_PROBE_RESP,
	VIRTIO_WIFI_BEACON_IES_ASSOC_RESP,
	VIRTIO_WIFI_BEACON_PROBE_RESP,
	VIRTIO_WIFI_BEACON_SSID,
	VIRTIO_WIFI_BEACON_AUTH_TYPE,
};
\end{lstlisting}

% TODO beacon fields

After an access point has started operation, the device can respond with asynchronous VIRTIO_WIFI_EVENT_CODE_NEW_STATION or VIRTIO_WIFI_EVENT_CODE_DELETE_STATION events.

The VIRTIO_WIFI_EVENT_CODE_NEW_STATION and VIRTIO_WIFI_EVENT_CODE_DELETE_STATION events have the type:

\begin{lstlisting}
struct virtio_wifi_station_change {
	struct virtio_wifi_event hdr;
	__u8 mac[ETH_ALEN];
	__u8 padding[2];
	__le64 ie_len;
};
\end{lstlisting}

% TODO explain

\subsubsection{Registering Reception of Specific 802.11 Frames}\label{sec:Device Types / 802.11 Device / Device Operation / Registering Reception of Specific 802.11 Frames}

A driver can request to receive specific 802.11 management frames with the command VIRTIO_WIFI_CMD_REGISTER_FRAME.

\begin{lstlisting}
struct virtio_wifi_register_frame {
	struct virtio_wifi_hdr hdr;
	__le16 ftype;
	__le16 stype;
};
\end{lstlisting}

\begin{description}
\item[\field{ftype}] contains the frame type.
\item[\field{stype}] contains the frame subtype.
\end{description}

The device responds with a \field{struct virtio_wifi_hdr} indicating success or failure.

\subsubsection{Key Management}\label{sec:Device Types / 802.11 Device / Device Operation / Key Management}

Common key structure:

\begin{lstlisting}
struct virtio_wifi_key {
	struct virtio_wifi_hdr hdr;
	__u8 mac[ETH_ALEN];
	__le16 vlan_id;
	__le32 cipher;
	/* enum nl80211_key_mode */
	__le32 mode;
#define VIRTIO_WIFI_KEY_F_DEFAULT_UNI (1 << 0)
#define VIRTIO_WIFI_KEY_F_DEFAULT_MULTI (1 << 1)
#define VIRTIO_WIFI_KEY_F_DEFAULT_MGMT (1 << 2)
#define VIRTIO_WIFI_KEY_F_DEFAULT (1 << 3)
#define VIRTIO_WIFI_KEY_F_IDX (1 << 4)
#define VIRTIO_WIFI_KEY_F_TYPE (1 << 5)
	__le16 flags;
	__u8 pairwise;
	__u8 key_index;
};

enum virtio_wifi_key_field {
       VIRTIO_WIFI_KEY_DATA = 0,
       VIRTIO_WIFI_KEY_SEQ,
};
\end{lstlisting}

\begin{description}
\item[\field{flags}] TODO explain flags
\item[\field{cipher}] selects the cipher suite.
\item[\field{VIRTIO_WIFI_KEY_DATA}] contains the key material.
\item[\field{VIRTIO_WIFI_KEY_SEQ}] contains the key sequence counter and is only used with VIRTIO_WIFI_CMD_GET_KEY.
\end{description}

\paragraph{Querying Keys}\label{sec:Device Types / 802.11 Device / Device Operation / Key Management / Querying keys}

To retrieve a specific key, the driver submits a \field{struct virtio_wifi_key} along with key and sequence bytes to the controlq with \field{code} set to VIRTIO_WIFI_CMD_GET_KEY.

The device responds with a \field{struct virtio_wifi_key} containing the requested key.

\paragraph{Adding Keys}\label{sec:Device Types / 802.11 Device / Device Operation / Key Management / Adding keys}

To add a key, the driver submits a \field{struct virtio_wifi_key} along with key and sequence bytes to the controlq with \field{code} set to VIRTIO_WIFI_CMD_ADD_KEY.

The device responds with a \field{struct virtio_wifi_hdr} indicating success or failure.

\paragraph{Deleting Keys}\label{sec:Device Types / 802.11 Device / Device Operation / Key Management / Deleting keys}

To delete a key, the driver submits a \field{struct virtio_wifi_key} (setting only fields \field{key_index} and \field{pairwise}) to the controlq with \field{code} set to VIRTIO_WIFI_CMD_DEL_KEY.

The device responds with a \field{struct virtio_wifi_hdr} indicating success or failure.

\paragraph{Modifying Keys}\label{sec:Device Types / 802.11 Device / Device Operation / Key Management / Modifying keys}

To modify key properties \field{unicast} and \field{multicast}, the driver submits a \field{struct virtio_wifi_key} (setting only fields \field{key_index}, \field{unicast} and \field{multicast}) to the controlq with \field{code} set to VIRTIO_WIFI_CMD_SET_KEY.

The device responds with a \field{struct virtio_wifi_hdr} indicating success or failure.

\subsubsection{Sending Packets (TX)}\label{sec:Device Types / 802.11 Device / Device Operation / Packet TX}

The driver writes a \field{struct virtio_wifi_hdr} with the \field{index} set to the interface of the packet, followed by the packet data.

\field{code} can be either VIRTIO_WIFI_802_3 or VIRTIO_WIFI_802_11.

\subsubsection{Receiving Packets (RX)}\label{sec:Device Types / 802.11 Device / Device Operation / Packet RX}

The driver reads a \field{struct virtio_wifi_hdr} with the \field{index} set to the interface of the packet, followed by the packet data.

\field{code} can be either VIRTIO_WIFI_802_3 or VIRTIO_WIFI_802_11.
